% Drafted against audit/outline.md §5. Three figures in that block were stale
% and are corrected here against the current suite and data: sizing error is
% 4.26e-14 percent rather than 6.6e-14, the Kupiec empirical size is 4.3% at
% the seed the test fixes rather than 4.8-5.8%, and the suite is 107 tests
% rather than 88. The outline records the corrections.
%
% Berkowitz and O'Brien (2002) is not cited: paywalled, never read, recorded
% as such in the source register.
%
% Evidence: data/processed/calibration_error_by_year.csv and the test suite,
% hash-recorded in audit/evidence_freeze.csv at commit 1e6bdd74.

\section{Validation}\label{sec:validation}

A model that runs is not a model that is right. Every check in this section
compares the engine against something outside it --- a hand calculation, an
independent implementation, a published boundary, an algebraic identity --- so
that agreement is evidence rather than self-report. The suite is 107 tests.
What none of it establishes is the resemblance between the proxy book and
Goldman's, and the last subsection says so plainly.

\subsection{The pricer against hand calculation}

The bond pricer is the one component where an error would be silent: a
mispriced bond produces a plausible VaR rather than an exception. A ten-year
bond with a 5\% semiannual coupon valued at a 6\% yield prices at
92.5612626, which the implementation reproduces to seven decimals, and it
trades below par because its coupon is a point under its yield. A zero-coupon
bond reproduces $100/1.03^{10}$ to twelve decimals at a 3\% rate. A par bond
prices at exactly 100 at its own yield, which is the identity that catches a
coupon-frequency error.

Two further properties are checked because
Equation~\eqref{eq:revaluation} depends on them. The analytic modified duration
agrees with a numerical derivative of the price function, and convexity has the
correct sign in both directions: a yield fall gains more than duration predicts
and a yield rise loses less. That second property is what
Section~\ref{sec:assumptions} quantifies when it prices the rate sleeve both
ways.

\subsection{The sizing identity}

Section~\ref{sec:method} claims the sleeve calibration is exact rather than
approximate, on the grounds that VaR is positively homogeneous in notional.
Table~\ref{tab:sizing} checks the claim instead of assuming it. Across all 40
sleeve-years the largest absolute sizing error is
$4.26 \times 10^{-14}$ percent of target, which is floating-point noise.

% Generated by build/make_tables.py from calibration_error_by_year.csv. Do not edit.

\begin{table}[tbp]
\centering
\small
\caption{Sizing check: the disclosed sleeve targets, and the largest error the calibration leaves behind}
\label{tab:sizing}
\begin{tabular}{@{}lrrrrr@{}}
\toprule
 & \multicolumn{4}{c}{Disclosed target (\$mm)} & Largest error \\
\cmidrule(lr){2-5}
Year & Rates & Equity & Currency & Commodity & (\% of target) \\
\midrule
2016 & 45 & 25 & 21 & 17 & 4.26e-14 \\
2017 & 40 & 24 & 12 & 13 & 1.78e-14 \\
2018 & 46 & 31 & 14 & 11 & 1.61e-14 \\
2019 & 46 & 27 & 11 & 12 & 3.23e-14 \\
2020 & 71 & 55 & 23 & 20 & 3.09e-14 \\
2021 & 60 & 43 & 13 & 25 & 3.30e-14 \\
2022 & 97 & 33 & 32 & 47 & 3.02e-14 \\
2023 & 96 & 29 & 24 & 19 & 2.96e-14 \\
2024 & 81 & 37 & 26 & 19 & 3.84e-14 \\
2025 & 70 & 47 & 24 & 17 & 3.02e-14 \\
\bottomrule
\end{tabular}
\notetext{Targets are Goldman's disclosed average daily VaR by category, as first reported. The final column is the largest absolute sizing error across the four sleeves in that year, as a percentage of target. Across all 40 sleeve-years the largest error is 4.26e-14\%, which is floating-point noise: the calibration is exact, as Equation~\eqref{eq:homogeneity} requires, and this table verifies that claim rather than reporting a convergence result.}
\end{table}


Two things make this worth a table rather than a footnote. The first is that an
inexact calibration would quietly contaminate the comparison: if summed
category VaR only approximately matched the disclosures, part of any gap in the
diversification effect would be a sizing artefact rather than a correlation
result. The second is that the check is independent of the mechanism. The
engine is not told that scaling is linear; it is run, and the achieved sleeve
VaR is measured and compared against target. Agreement to fourteen decimal
places is then evidence for the homogeneity property, not a restatement of it.

A test in the suite makes the same point directly, by running the whole engine
on a book and on the same book scaled by an arbitrary factor and confirming
that every VaR scales by exactly that factor.

\subsection{The backtest statistics against published values}

The coverage tests are standard and their implementations are checked against
external values rather than against each other.

The \citet{kupiec1995} proportion-of-failures statistic reproduces an
independent implementation on two published cases. With 1{,}043 observations and
57 failures at 95\%, the published statistic is 0.46147 with a $p$-value of
0.49694, and this implementation returns 0.4614664, agreeing to the five decimal
places the published figure carries. With 17 failures at 99\% the published
statistic is 3.51180 and this implementation returns 3.5118129, a difference of
$1.3 \times 10^{-5}$; the published figure is not the correctly rounded
five-decimal value, which is 3.51181, so the agreement is to within that
tolerance rather than digit for digit. The $p$-value, 0.060933, matches
exactly. Its analytic limits hold as well. Zero exceptions give
$-2T\ln(1-p)$, and a failure rate exactly equal to coverage gives zero. A
Monte Carlo experiment under the null puts the empirical size at 4.3\% against
a nominal 5\% at the seed the test fixes, and between 4.4\% and 5.7\% across
three further seeds. The statistic is discrete, so its distribution cannot
match a continuous chi-square exactly; size at the nominal level is the
property that matters and the one that is tested.

The \citet{christoffersen1998} independence and conditional coverage statistics
are checked against their own analytic cases, including the degenerate samples
where a transition count is zero.

The Basel traffic-light boundaries deserve a note because they are easy to get
wrong by borrowing them. They are rederived here from the binomial
distribution, not hard-coded, and the derivation reproduces the published
boundaries of five and ten for a 99\% model over 250 trading days
\citep{bcbs1996backtesting}. The reason the derivation matters is that those
numbers are specific to 99\%. A correctly calibrated 95\% model is expected to
break about twelve times in 250 days, and the rederivation puts its yellow and
red thresholds at 18 and 27; carrying the published 99\% boundaries across
unchanged would place a well-specified 95\% model permanently in the red zone.
Every zone reported in this paper is computed at the confidence level it
belongs to.

\subsection{The look-ahead probe}

Look-ahead is the failure mode that announces itself as success. Every other
error in a VaR engine produces an implausible number; using tomorrow's data
produces an unusually good backtest, which is easy to mistake for a working
model.

The probe is therefore adversarial. Every factor move after a cut-off date is
multiplied by 25, and the engine is re-run on the tampered panel. VaR stamped
on or before the cut-off must be identical to what it was, because it is only
permitted to see scenarios up to its own date. The realised profit-and-loss
column is excluded from the comparison, since that column is supposed to
change: it carries the next day's outcome, which is exactly the quantity the
tampering alters.

A companion test pins the off-by-one that would make the backtest meaningless,
confirming that the realised result recorded against date $t$ is the book
repriced on the move of $t+1$ and not of $t$.

\subsection{The subadditivity guard}

One check needs explaining rather than asserting, because its natural form
would be wrong.

The diversification effect in Equation~\eqref{eq:diversification} is almost
always negative, and an engine that produced a positive one would normally be
signalling a sign error. But VaR is not a subadditive risk measure.
\citet{artzner1999coherent} define a coherent measure by four axioms:
translation invariance, subadditivity, positive homogeneity and monotonicity.
They show that VaR satisfies the first, third and fourth while failing the
second. In their words, the use of VaR ``does not encourage and, indeed,
sometimes prohibits diversification.'' On a sufficiently awkward joint tail, a
whole book really can show a larger loss quantile than the sum of its parts.
That is a property of the measure rather than a defect in the code, and it is
the reason expected shortfall exists.

So the guard is a tolerance, not an assertion. The engine records days on which
the effect is positive and refuses the run only when they exceed one day in a
hundred, which is the threshold separating the genuine phenomenon from a sign
error. A test confirms the guard behaves that way in both directions: it
tolerates a few such days and rejects many.

On every configuration reported in this paper the count is zero, 0 of 2{,}293
days on both the fixed and the recalibrated book. That is not an accident of
this data. \citet{artzner1999coherent} note in a remark that under joint
normality the quantiles do satisfy subadditivity as long as exceedance
probabilities stay below one half, and a four-sleeve book of broad market
factors at a 5\% tail is close enough to that case for the failure not to bite.
The guard exists for the configurations where it might, not because it ever
fired.

\subsection{The equal-weight collapse}

One check exists specifically to make Section~\ref{sec:assumptions}
interpretable. The age-weighted estimator of
Equation~\eqref{eq:ageweights} and the ordinary percentile are different pieces
of code, and if they disagreed at $\lambda = 1$ then any difference between a
weighted and an unweighted run would mix the weighting with a change of
estimator.

They do not disagree. The plotting position in the weighted reading is chosen so
that equal weights reduce exactly to the position the unweighted percentile
interpolates against, and the discrete readings coincide likewise. The suite
pins the identity to machine precision across four sample sizes and four
quantiles, in both conventions. Whatever Section~\ref{sec:assumptions} reports
about age weighting is therefore the weighting.

\subsection{What the engine does not validate}

Everything above establishes internal correctness: the engine computes what
Section~\ref{sec:method} says it computes, on the data
Section~\ref{sec:sample} describes. Within that, two things are checked by the
build rather than by the suite and should be named: the alignment of the factor
panel to a common trading calendar, and the staleness and sign checks on the
raw series. Those run on every rebuild and halt it on failure, but they are not
unit tests and no test asserts that the calendar intersection is correct beyond
what those checks cover. None of it establishes that the proxy book
resembles Goldman's trading book, and no test could, because the positions are
not disclosed and never will be.

This matters for how the results should be read. A validated engine on a proxy
book can show that a particular structure of market risk, exposed to public
factor data, does or does not generate a particular diversification effect.
It cannot show that Goldman's book has that structure. The level agreement in
Section~\ref{sec:results} is therefore evidence about the proxy, and the
steadiness gap is evidence about what a static proxy cannot do, not a
measurement of anything inside the firm.
